\documentclass[11pt]{article}
\usepackage{icmlko}
\begin{document}

\title{여섯 점처럼 보였지만 두 점이었다 --- 내가 만든 변이가 내 설명을
안 받쳐 준다}
\author{월드모델 실험실}
\date{노트 467 · 2026년 8월 2일}
\maketitle

\begin{abstract}
노트 466 이 후보 규칙을 하나 세웠다 --- \emph{모델이 그 축을 축값에
단조롭게 쓰면 두고, 거의 안 쓰면 가린다.} \textbf{점 둘로 세운 규칙}이라
문턱도 못 정했고 우연일 수도 있었다. 그래서 변이를 직접 만들었다 ---
축을 하나씩 빼서 짝마다 팔 셋, \textbf{모두 여섯 팔}.

\begin{center}
\begin{tabular}{lrrr}
\hline
팔 & 남는 축 & 쓰임$\leftrightarrow$축값 & 가림 효과 \\
\hline
KR 만화 & $4$ & $-0.604$ & $-0.0584$ \\
KR 만화 & $3$ & $-0.554$ & $-0.0955$ \\
KR 만화 & $2$ & $-0.270$ & $-0.0855$ \\
비게임 앱 & $5$ & $-0.026$ & $+0.0435$ \\
비게임 앱 & $4$ & $+0.350$ & $+0.0830$ \\
비게임 앱 & $3$ & $+0.238$ & $+0.0761$ \\
\hline
\end{tabular}
\end{center}

여섯 점을 한 줄에 놓으면 \textbf{rho $+0.829$}(p $0.042$)로 그럴듯하다.
\textbf{반증이다.}

\begin{itemize}
\item \textbf{짝 안에서 순서가 안 맞는다} --- KR 세 팔의 상관이
$\mathbf{-0.500}$ 으로 \emph{거꾸로}다(앱만 $+1.000$).
\item \textbf{문턱이 안 선다} --- KR $2$축은 쓰임이 $-0.270$ 으로 내가
잡은 문턱 $-0.3$ 을 \emph{넘는데} 효과는 $-0.0855$ 로 음수다.
\end{itemize}

\textbf{상관 $+0.829$ 를 만드는 것은 짝 사이 간격 하나다} --- 그리고 그
간격이야말로 내가 \emph{설명하려던} 것이다. 설명이 설명 대상을 다시
가리키고 있었다.
\end{abstract}

\begin{figure}[t]
\centering
\includegraphics[width=.88\textwidth]{figs/f1.pdf}
\caption{여섯 팔. 회귀선은 그럴듯하지만 두 무리가 각자 뭉쳐 있고, KR
무리 안에서는 기울기가 반대다.}
\end{figure}

\section{등록할 때 적어 둔 주의가 맞았다}

사전 등록에 이렇게 적었다 --- \emph{``여섯 팔이 독립이 아니다(같은 짝에서
나온 셋씩) --- 그래서 짝 안에서도 순서가 맞는지 따로 본다.''} 그 조항이
없었다면 rho $+0.829$ · p $0.042$ 를 보고 규칙이 섰다고 적었을 것이다.

\textbf{변이를 만드는 것과 변이를 얻는 것은 다르다.} 축을 빼서 만든
팔들은 \emph{같은 수집의 같은 행}이고, 바뀐 것은 모델에게 보여 주는
축뿐이다. 그래서 이 팔들은 \textbf{수집 사이의 차이를 설명할 힘이 없다}
--- 애초에 수집을 안 바꿨으니까.

\begin{quote}
\textbf{조항} --- 무리 사이 차이를 설명하려고 무리 \emph{안}에 변이를
만들었으면, \textbf{무리 안 관계를 따로 적는다.} 뭉친 점들을 한 줄에
놓고 잰 상관은 설명하려던 차이를 다시 재는 것일 수 있다.
\end{quote}

\section{그러면 무엇이 남나}

\textbf{쓰임$\leftrightarrow$축값이 두 수집을 가른다는 것은 여전히
사실이다}(KR 셋 다 음수 $-0.60 \sim -0.27$, 앱 셋 다 $-0.03$ 이상).
하지만 그것이 \emph{원인}인지 \emph{또 하나의 이름표}인지는 이 실험으로
못 정한다. 두 짝을 가르는 어떤 값이든 두 짝을 가른다.

노트 466 의 후보 규칙은 \textbf{가설로 남기되 근거를 내린다.} 검증은
\textbf{새 수집}으로만 된다 --- 라벨 보기 전에 쓰임을 재서 예보하고,
그 다음 라벨로 채점하는 것. 그때까지 \texttt{OOC\_DROP} 은 노트 463 대로
\textbf{미결}이다.

\section*{남은 것}

\begin{itemize}
\item \textbf{짝 안 변이를 진짜로 만들려면} 축이 아니라 \emph{행}을
갈라야 한다 --- 노트 465 의 겹 셋이 그것인데, 겹은 가림 효과가
$0.02$ 안에서만 움직여 쓰임도 거의 안 움직일 것이다. 더 센 갈래(연도 ·
갈래)로 나누면 안쪽 변이가 생기나.
\item \textbf{셋째 수집} --- 이 물음의 모든 길이 거기로 모인다.
\item 노트 466 이 남긴 다른 물음(다른 넷도 수집마다 쓰임이 갈리나)은
그대로 열려 있다.
\end{itemize}

\end{document}
